\documentclass{amsart}
\begin{document}

The aim is to show that strange Besov norm defined by the formula 
$$\|f\|_{B^{p}_{Dunkl}(\mathbb{R}^N)}:=\left(\int_{\mathbb{R}^{N}}\int_{\mathbb{R}^{N}}\left(\frac{d(x,y)}{\|x-y\|}\right)^{p}\frac{|f(x)-f(y)|^p}{w(B(x,d(x,y)))^2}w(x)w(y)dxdy\right)^{1/p}$$
is controlled by the usual Sobolev norm $\|f\|_{\dot{W}^{\frac{N}{p},p}(\mathbb{R}^N)}.$ In this file, we only deal with the following example.

Weyl group $\mathbb{Z}_2^2$ acting on $\mathbb{R}^2$ is equipped with multiplicity function $k(2^{\frac12},0)=a$ and $k(0,2^{\frac12})=b.$ We have
$$w(x)=|x_1|^{2a}|x_2|^{2b}.$$
We have
$$d(x,y)\sim \big||x_1|-|y_1|\big|+\big||x_2|-|y_2|\big|.$$
Thus, for 
$$w(B(x,d(x,y)))\sim \big||x_1|-|y_1|\big|+\big||x_2|-|y_2|\big|^2\times \big(|x_1|+|y_1|+\big||x_2|-|y_2|\big|\big)^{2a}\times$$
$$\times \big(|x_2|+|y_2|+\big||x_1|-|y_1|\big|\big)^{2b}.$$

Consider, for example, the first quadrant for $x$ and the second quadrant for $y.$ We have
$$\int_{\mathbb{R}^2_+}\int_{\mathbb{R}^2_+}\frac{(|x_1-y_1|+|x_2-y_2|)^{p-4}}{(x_1+y_1+|x_2-y_2|)^{p+4a}}\frac{|f(x)-f(s_1(y))|^p}{(x_2+y_2+|x_1-y_1|)^{4b}}(x_1y_1)^{2a}(x_2y_2)^{2b}dxdy.$$
We already know that
$$\int_{\mathbb{R}^2_+}\int_{\mathbb{R}^2_+}\frac{(|x_1-y_1|+|x_2-y_2|)^{p-4}}{(x_1+y_1+|x_2-y_2|)^{p+4a}}\frac{|f(x)-f(y)|^p}{(x_2+y_2+|x_1-y_1|)^{4b}}(x_1y_1)^{2a}(x_2y_2)^{2b}dxdy$$
is controlled by the Besov norm in the first quadrant. We therefore, need to control
$$\int_{\mathbb{R}^2_+}\int_{\mathbb{R}^2_+}\frac{(|x_1-y_1|+|x_2-y_2|)^{p-4}}{(x_1+y_1+|x_2-y_2|)^{p+4a}}\frac{|f(y)-f(s_1(y))|^p}{(x_2+y_2+|x_1-y_1|)^{4b}}(x_1y_1)^{2a}(x_2y_2)^{2b}dxdy=$$
$$=\int_{\mathbb{R}^2_+}\Big(\int_{\mathbb{R}^2_+}\frac{(|x_1-y_1|+|x_2-y_2|)^{p-4}(x_1y_1)^{2a}(x_2y_2)^{2b}dx}{(x_1+y_1+|x_2-y_2|)^{p+4a}(x_2+y_2+|x_1-y_1|)^{4b}}\Big)|f(y)-f(s_1(y))|^pdy.$$

Similar expression for the usual Sobolev norm is
$$\int_{\mathbb{R}^2_+}\Big(\int_{\mathbb{R}^2_+}\frac{dx}{(x_1+y_1+|x_2-y_2|)^4}\Big)|f(y)-f(s_1(y))|^pdy\approx y_1^{-2}.$$

So, we need to establish
$$\int_{\mathbb{R}^2_+}\frac{(|x_1-y_1|+|x_2-y_2|)^{p-4}(x_1y_1)^{2a}(x_2y_2)^{2b}dx}{(x_1+y_1+|x_2-y_2|)^{p+4a}(x_2+y_2+|x_1-y_1|)^{4b}}\leq y_1^{-2},\quad y\in\mathbb{R}^2_+.$$

Since
$$(x_1y_1)^{2a}\leq (x_1+y_1+|x_2-y_2|)^{p+4a},\quad (x_2y_2)^{2b}\leq (x_2+y_2+|x_1-y_1|)^{4b},$$
it follows that one needs to establish
$$\int_{\mathbb{R}^2_+}\frac{(|x_1-y_1|+|x_2-y_2|)^{p-4}dx}{(x_1+y_1+|x_2-y_2|)^p}\leq y_1^{-2},\quad y\in\mathbb{R}^2_+.$$

We need to consider $4$ regions: (A) $|x_1-y_1|\geq\frac12y_1$ and $|x_2-y_2|\geq\frac12y_2$ (B) $|x_1-y_1|\geq\frac12y_1$ and $|x_2-y_2|\leq\frac12y_2$ (C) $|x_1-y_1|\leq\frac12y_1$ and $|x_2-y_2|\geq\frac12y_2$ (D) $|x_1-y_1|\leq\frac12y_1$ and $|x_2-y_2|\leq\frac12y_2.$

{\bf Case A:} we have $|x_1-y_1|\sim x_1+y_1$ and $|x_2-y_2|\sim x_2+y_2.$ We, therefore, need to check
$$\int_{\mathbb{R}^2_+}\frac{dx}{(x_1+y_1+x_2+y_2)^4}\leq y_1^{-2},\quad y\in\mathbb{R}^2_+.$$
The left hand side becomes maximal if $y_2=0.$ We, therefore, need to check
$$\int_{\mathbb{R}^2_+}\frac{dx}{(x_1+y_1+x_2)^4}\leq y_1^{-2},\quad y\in\mathbb{R}^2_+.$$
By dilation invariance, it is equivalent to
$$\int_{\mathbb{R}^2_+}\frac{dx}{(x_1+1+x_2)^4}<\infty.$$
This is obvious.

{\bf Case B:} we have $|x_1-y_1|\sim x_1+y_1.$ We, therefore, need to check
$$\int_{\mathbb{R}_+\times(\frac12y_2,\frac32y_2)}\frac{dx}{(x_1+y_1+|x_2-y_2|)^4}\leq y_1^{-2},\quad y\in\mathbb{R}^2_+.$$
Using new variable $u=x_2-y_2,$ we equivalently rewrite it as
$$\int_{\mathbb{R}_+\times(-\frac12y_2,\frac12y_2)}\frac{dx_1du}{(x_1+y_1+|u|)^4}\leq y_1^{-2},\quad y\in\mathbb{R}^2_+.$$
The left hand side grows with $y_2.$ It, therefore, suffices to prove
$$\int_{\mathbb{R}_+\times\mathbb{R}}\frac{dx_1du}{(x_1+y_1+|u|)^4}\leq y_1^{-2},\quad y\in\mathbb{R}^2_+.$$
By dilation invariance, it is equivalent to
$$\int_{\mathbb{R}_+\times\mathbb{R}}\frac{dx_1du}{(x_1+1+|u|)^4}<\infty.$$
This is obvious.

{\bf Case C:} we have $x_1+y_1\sim y_1$ and $|x_2-y_2|\sim x_2+y_2.$ We, therefore, need to check
$$\int_{(\frac12y_1,\frac32y_1)\times\mathbb{R}_+}\frac{(|x_1-y_1|+x_2+y_2)^{p-4}dx}{(y_1+x_2+y_2)^p}\leq y_1^{-2},\quad y\in\mathbb{R}^2_+.$$
Using new variable $u=x_1-y_1,$ we equivalently rewrite it as
$$\int_{(-\frac12y_1,\frac12y_1)\times\mathbb{R}_+}\frac{(|u|+x_2+y_2)^{p-4}dudx_2}{(y_1+x_2+y_2)^p}\leq y_1^{-2},\quad y\in\mathbb{R}^2_+.$$
Substituting $u=y_1z_1,$ $x_2=y_1z_2$ and denoting $t=\frac{y_2}{y_1},$ we equivalently rewrite it as
$$\int_{(-\frac12,\frac12)\times\mathbb{R}_+}\frac{(|z_1|+z_2+t)^{p-4}dz}{(1+z_2+t)^p}\leq 1,\quad t>0.$$
It suffices to prove
$$\int_{(0,1)\times\mathbb{R}_+}\frac{(z_1+z_2+t)^{p-4}dz}{(1+z_2+t)^p}\leq 1,\quad t>0.$$
This is equivalent to a combination of inequalities
$$\int_{(0,1)\times(0,1)}\frac{(z_1+z_2+t)^{p-4}dz}{(1+z_2+t)^p}\leq 1,\quad t>0.$$
$$\int_{(0,1)\times(1,\infty)}\frac{(z_1+z_2+t)^{p-4}dz}{(1+z_2+t)^p}\leq 1,\quad t>0.$$
In other words,
$$\int_{(0,1)\times(0,1)}\frac{(z_1+z_2+t)^{p-4}dz}{(1+t)^p}\leq 1,\quad t>0.$$
$$\int_{(0,1)\times(1,\infty)}(z_2+t)^{-4}dz\leq 1,\quad t>0.$$
In other words,
$$\int_{(0,1)\times(0,1)}(z_1+z_2+t)^{p-4}dz\leq (t+1)^p,\quad t>0.$$
$$\int_{(0,1)\times(1,\infty)}z_2^{-4}dz<\infty.$$
The second inequality is obvious. The first inequality is equivalent to
$$\int_{|z|<1}(|z|+t)^{p-4}dz\leq (t+1)^p,\quad t>0.$$
In other words, we need to check
$$\int_0^1(r+t)^{p-4}rdr\leq (t+1)^p,\quad t>0.$$

If $p<4,$ then
$$\int_0^1(r+t)^{p-4}rdr\leq \int_0^1r^{p-4}rdr=\frac1{p-2}\leq (t+1)^p,\quad t>0.$$
If $p\geq 4,$ then
$$\int_0^1(r+t)^{p-4}rdr\leq \int_0^1r^{p-3}dr+\int_0^1t^{p-4}rdr=\frac1{p-2}+\frac12t^{p-4}\leq (t+1)^p,\quad t>0.$$
Hence, we verified the inequality in this (the most complicated) case.

{\bf Case D:} we have $x_1+y_1\sim y_1$ and $x_2+y_2\sim y_2.$ We, therefore, need to check 
$$\int_{(\frac12y_1,\frac32y_1)\times(\frac12y_2,\frac32y_2)}\frac{(|x_1-y_1|+|x_2-y_2|)^{p-4}dx}{(y_1+|x_2-y_2|)^p}\leq y_1^{-2},\quad y\in\mathbb{R}^2_+.$$
Using new variables $u=x_1-y_1$ and $v=x_2-y_2,$ we equivalently rewrite it as
$$\int_{(-\frac12y_1,\frac12y_1)\times(-\frac12y_2,\frac12y_2)}\frac{(|u|+|v|)^{p-4}dudv}{(y_1+|v|)^p}\leq y_1^{-2},\quad y\in\mathbb{R}^2_+.$$
The left hand side grows with $y_2.$ It, therefore, suffices to prove
$$\int_{(-\frac12y_1,\frac12y_1)\times\mathbb{R}}\frac{(|u|+|v|)^{p-4}dudv}{(y_1+|v|)^p}\leq y_1^{-2},\quad y\in\mathbb{R}^2_+.$$
By dilation invariance, this is equivalent to
$$\int_{(-\frac12,\frac12)\times\mathbb{R}}\frac{(|u|+|v|)^{p-4}dudv}{(1+|v|)^p}<\infty.$$
As we assume $p>2,$ this is obvious.



So, we proved the estimate in each particular case.
\end{document}